% Generated from locale/tr/content/computability/computability-theory/coding-computations.tex; do not hand-edit.


\olfileid[tr]{cmp}{thy}{cod}
\olsection{Hesapların Kodlanması}

Her hesaplama modelinde aşağıdakileri yapmak mümkündür:
\begin{enumerate}
\item Hesaplanabilir fonksiyonların \emph{tanımlarını} sistematik biçimde
  betimlemek. Örneğin Turing makinesi belirtimleri, özyinelemeli tanımlar ya
  da bir programlama dilindeki programlar bu tanımları sağlar.
\item Belirli bir tanımla verilen fonksiyonun belirli bir girdideki hesabının
  eksiksiz kaydını betimlemek. Örneğin bir Turing makinesi hesabı, her hesap
  adımındaki konfigürasyonların, yani makinenin durumu ile şerit içeriğinin
  dizisiyle betimlenebilir.
\item Bir hesabın varsayılan kaydının, belirli bir tanıma sahip hesaplanabilir
  fonksiyonun belirli bir girdide, fonksiyonun tanımlı ve hesabın durduğu
  durumda, nasıl hesaplandığının gerçekten kaydı olup olmadığını sınamak.
\item Bir hesabın tam kaydına ilişkin böyle bir betimden, fonksiyonun belirli
  girdideki değerini çıkarmak. Örneğin duran bir Turing makinesi hesabının
  en son adımındaki şerit içeriği değerdir.
\end{enumerate}

Kodlama kullanılarak hesaplanabilir bir fonksiyonun her betimine sayısal bir
\emph{indis}, yönergeler bu indisten hesaplanabilir biçimde geri elde
edilebilecek şekilde atanabilir. Benzer biçimde bir hesabın tam kaydı da tek
bir sayıyla kodlanabilir. Böylece ortaya çıkan ``$s$, indisi $e$ olan
fonksiyonun $x$ girdisindeki hesabının kaydını kodlar'' aritmetik bağıntısı
ile ``kodu $s$ olan hesap dizisinin çıktısı'' fonksiyonu hesaplanabilirdir;
hatta ilkel özyinelemelidir.

Bu temel olgu çok güçlüdür ve seçilen hesaplama modelinden bağımsız olarak
hesaplanabilirlik hakkında çarpıcı ve önemli birçok sonuç ispatlamamızı
sağlar.

